\documentclass[journal]{IEEEtran}
\usepackage{amsmath,amssymb,amsthm,mathtools}
\usepackage{cite}
\usepackage{graphicx}
\usepackage[hidelinks]{hyperref}

\setlength{\emergencystretch}{2em}

\newtheorem{lemma}{Lemma}
\newtheorem{theorem}{Theorem}
\newtheorem{corollary}{Corollary}

\newcommand{\RegSDAC}[1]{\mathrm{Reg}_{\scriptscriptstyle\mathrm{SDAC}}(#1)}
\newcommand{\RegLin}[1]{\mathrm{Reg}_{\mathrm{lin}}(#1)}

\title{Online Constrained Control of Storage Systems via Simplex Disturbance-Action
Policies}

\author{Kamiar~Asgari and Michael~J.~Neely%
\thanks{K.~Asgari and M.~J.~Neely are with the Department of Electrical and
Computer Engineering, University of Southern California, Los Angeles, CA 90089 USA
(e-mail: kamiaras@usc.edu; mjneely@usc.edu).}%
\thanks{Corresponding author: K.~Asgari.}%
}

\markboth{IEEE Transactions on Automatic Control}%
{Asgari and Neely: Online Constrained Control of Storage Systems via SDAC}

\begin{document}
\maketitle

\begin{abstract}
We study online control of a scalar system under adversarial arrivals (often called
disturbances in the online-control literature), convex cost functions, and
state-dependent action constraints. At each time, the controller must choose a
feasible action before observing the corresponding arrival and cost. We introduce
\emph{Simplex Disturbance-Action Control} (SDAC), a simplex-constrained
disturbance-action policy class that preserves per-step feasibility by construction.
An online controller updates SDAC parameters via online mirror descent (OMD) on the
sub-probability simplex, with a feasibility projection that enforces the
state-dependent action constraint. We prove that regret against the best fixed SDAC
policy in hindsight is sublinear in the horizon and logarithmic in the policy
dimension, implying no-regret relative to feasible linear state-feedback policies.
The setting is motivated by energy-harvesting battery management and similar storage
systems.
\end{abstract}

\begin{IEEEkeywords}
Online control, disturbance-action policies, storage systems, energy harvesting,
regret minimization, state constraints.
\end{IEEEkeywords}

\section{Introduction}
\label{sec:introduction}

\subsection{Motivation and Problem Setting}
\IEEEPARstart{M}{any} dynamic resource systems can be modeled by a storage state that
evolves over time and constrains the control action at every time. The controller
must choose an action subject to state-dependent constraints before observing the
corresponding arrival and cost.

This creates a constrained discrete-time control problem with state--control
coupling: aggressive early control actions can harm later performance, while overly
conservative actions can increase per-step cost. Our objective is to design a policy
that remains feasible at every time and competes with the best fixed policy in a rich
policy class—one expressive enough to capture many strong policies.

\subsection{Model and Information Structure}
We study a discrete-time storage system with state $x_t \geq 0$, where $x_t$ models
the storage level at the beginning of time $t$. At each time $t$, the controller
chooses a control action $u_t$ under the \emph{state constraints}
\[
0 \leq u_t \leq \alpha x_t,
\]
where $\alpha\in(0,1)$ is a fixed known state-retention factor (so a fraction
$1-\alpha$ of the state is lost each period). After the action is chosen, a
nonnegative adversarial arrival $w_t$ (a resource inflow) is revealed, and the state
evolves according to
\[
x_{t+1} = \alpha x_t - u_t + w_t.
\]
In the online-control literature this exogenous input is often called a
\emph{disturbance}; we reserve that term for that literature and for the names
Disturbance-Action Control (DAC) and Simplex Disturbance-Action Control (SDAC), and
otherwise refer to $w_t$ as an \emph{arrival}. Unlike typical process-noise
disturbances, which may take both positive and negative values, our model treats
$w_t$ specifically as a nonnegative arrival of resources into storage
($w_t\in[0,1]$). That modeling choice is essential for the storage interpretation
and for the feasibility properties of SDAC.

Unlike stochastic application models that rely on distributional assumptions, our
formulation is adversarial. Arrivals $w_t$ and costs $c_t(x, u)$ may vary
arbitrarily over time, and the controller chooses $u_t$ before observing either at
time $t$. The goal is low cumulative cost relative to the best fixed policy in a
later-defined policy class, while both our algorithm and the benchmark policy
maintain feasibility throughout.
 
\subsection{Application Domains}
This abstraction captures a broad set of domains:
\begin{itemize}
    \item \textbf{Energy Harvesting and Battery Management:} $x_t$ is battery state
        of charge (retention factor $\alpha$ after self-discharge), $w_t$ is
        harvested energy, and $u_t$ is discharge or consumption bounded by available
        energy $\alpha x_t$.
    \item \textbf{Queueing Networks and Buffer Management:} $x_t$ is queue/buffer
        occupancy, $w_t$ is
        packet arrivals, and $u_t$ is service bounded by $\alpha x_t$.
    \item \textbf{Inventory Control:} $x_t$ is on-hand inventory (retention factor
        $\alpha$ after spoilage/deterioration), $w_t$ is supply or replenishment, and
        $u_t$ is sales bounded by $\alpha x_t$.
    \item \textbf{Investments and Bank Accounts:} $x_t$ is liquid capital/balance
        (retention factor $\alpha$ after fees), $w_t$ is incoming cash flow, and
        $u_t$ is capital deployment or withdrawal bounded by $\alpha x_t$.
    \item \textbf{Water Reservoir and Dam Management:} $x_t$ is stored water
        (retention factor $\alpha$ after evaporation/seepage), $w_t$ is inflow, and
        $u_t$ is controlled release bounded by $\alpha x_t$.
\end{itemize}
We take energy harvesting as the running interpretation; the other domains share the
same retention-coupled state dynamics and state-dependent action constraints.

\subsection{Related Work}

\paragraph{Online control}
A recent line of work in online learning studies online control, where a linear
dynamical system is operated under adversarial arrivals and costs. The objective
is low time-average cost relative to the best policy in hindsight, rather than
classical stabilization or worst-case guarantees. Foundational results for linear
systems and adversarial settings include \cite{agarwal2019online,
agarwal2019logarithmic, hazan2020nonstochasticBandit, simchowitz2020improper}; see
\cite{hazan_introduction_2025} for a modern overview.

A central idea in this literature is \emph{Disturbance-Action Control}
(DAC),\footnote{``Disturbance-Action Control'' is the standard name in the
online-control literature. In this paper the underlying signal is the nonnegative
arrival $w_t$; we retain the DAC/SDAC names for continuity with that literature.}
a policy parameterization in which control actions are expressed as affine or linear
functions of the current state and a finite history of past arrivals. This
representation yields tractable surrogates for policy optimization and enables
efficient no-regret updates in adversarial environments. Extensions include partial
observability \cite{lale2020logarithmic} and bandit feedback
\cite{hazan2020nonstochasticBandit}. SDAC adapts this template to state-constrained
resource systems, where feasibility is coupled over time through the dynamics.

Online control provides a strong foundation for our problem class: it accommodates
arbitrary time-varying arrivals and unknown general convex cost functions, and the
resulting algorithms come with robust guarantees and tractability. However, for
state-coupled storage dynamics, unconstrained DAC does not enforce per-step
feasibility by construction. Our framework fills this gap by replacing DAC with SDAC,
a simplex-constrained DAC-style family that preserves the feasibility defined by
\eqref{eq:model-feasibility} while retaining the desirable characteristics.

\paragraph{Energy Harvesting}
A closely related work is \cite{lin2025optimal}, which studies an online
energy-harvesting allocation problem with dynamics similar to ours. That work
benchmarks against the best sequence of decisions in hindsight, but it differs in
several important ways: \cite{lin2025optimal} uses competitive ratio rather than
regret, adopts an action-only cost model rather than general convex state-action
costs, which means in particular that it cannot optimize a function of the battery
size, and it has a less adversarial information pattern because the algorithm
observes the current cost before selecting the action, whereas in our model the
control action is chosen before the cost is revealed.

Prior work on energy harvesting includes \cite{Yu2019}, which studies online
allocation with stochastic costs, and \cite{10.1145/3428337}, which improves over
\cite{Yu2019} by achieving $O(\sqrt{T})$ regret with finite battery and adversarial
costs. Both works consider action-only costs (not state-dependent), i.i.d.
arrivals (not adversarial), and benchmark against fixed allocation amounts rather
than fixed policies. In contrast, our framework jointly optimizes over state-action
costs, adversarial arrivals, and stronger benchmarks, requiring new
algorithmic techniques.

\paragraph{Queueing and Network Optimization}
An alternative, though mathematically related, perspective comes from queueing
control in computer networks and telecommunications. In that framework, one tracks a
queue with stochastic arrivals and makes real-time service decisions, typically
aiming to minimize long-run average service cost while keeping queue lengths stable.
This literature is extensive (see, e.g.,
\cite{tassiulas1990stability,tassiulas1993dynamic,mckeown1999achieving,neely2010stochastic}),
but its setting differs from ours in several important respects. Classical queueing
work usually assumes stochastic arrivals and focuses on simple, known cost functions
that depend only on the service rate, while queue length is kept small. These methods
achieve trade-offs between queue size and time-average service cost. By contrast, our
model permits fully adversarial sequences of arrivals and cost functions and allows
general convex costs that depend on both the queue length (state) and the service
rate (control action). Thus our approach is more general in cost structure and
exogenous-arrival model. The queueing literature provides useful motivation and
techniques for state-coupled dynamics, but does not directly address the online
optimization and regret-minimization objectives we study.

Moreover, in most practical queueing applications the natural retention factor is
$\alpha=1$ (no abandonment/leakage), so those settings are related dynamics but not
a direct instance of the $\alpha\in(0,1)$ model studied in this paper. The case
$\alpha\ge 1$ can be handled by a straightforward extension (Section~\ref{sec:extensions});
we do not develop it in here, as it adds technical overhead without
introducing new ideas.

Another related work is \cite{khojastepour2004delay}, which proposes and analyzes a
fixed controller that is a linear function of past arrivals. That controller
resembles our family of SDAC policies because it uses nonnegative filter weights that
sum to one (summing to at most one for SDAC). However, the controller in
\cite{khojastepour2004delay} is not updated online and is chosen to achieve bounded
service delay; by contrast, our framework focuses on online optimization of a general
cost and does not treat delay as a constraint.

\paragraph{Inventory Management}
A related work is \cite{yang2020online}, which studies a distinct model in which
demand (equivalent to allocation in our setup) is chosen adversarially, while the
controller chooses the supply (equivalent to arrivals in our model). The
objective is to minimize time-average purchasing cost while always meeting demand,
and the performance metric is competitive ratio rather than regret. Although
structurally similar to our setting, the control variable differs: we choose the
control action from the feasible state bound, whereas \cite{yang2020online} chooses
how much to procure. Another key difference is the cost model: our framework allows
general convex state-action costs, while \cite{yang2020online} uses a linear
cost depending only on the control action.

More closely related to ours is \cite{hihat2024online}, which considers an
inventory-control formulation with potentially nonlinear and nondifferentiable
state-transition functions, and studies base-stock policies. While that model is more
general in scope, its decision structure differs from ours: as in
\cite{yang2020online}, the controller chooses supply (inflow) and then observes
demand (allocation) through the transition, whereas our controller chooses the
control action and observes incoming supply (arrivals). Methodologically, our
setup is intentionally narrower and tailored to SDAC-based adversarial regret
analysis: feasibility constraints are enforced by policy design, yielding a simpler
linear state recursion, lower computational overhead, and cleaner analysis than
approaches operating through nonlinear nondifferentiable dynamics such as
\cite{hihat2024online}. The two frameworks are structurally related but target
different problem structures, control variables, and policy classes.

\paragraph{Water Reservoir and Dam Management}
Classical dam theory studies storage dynamics with random inflows and controlled
releases (see, for example, \cite{moran1954dams,prabhu1998stochastic}). Relative to
our adversarial online-control setting, this literature typically assumes stochastic
inputs and emphasizes stationary performance criteria under simple, known cost
structures. Here, it serves mainly as historical context for state-coupled control
under state constraints.

\subsection{Contributions}
The contributions are organized in two layers:
\begin{itemize}
    \item \textbf{Modeling and policy-class contribution.} We formulate adversarial
        online control of a storage system with retention-coupled state dynamics,
        state-dependent feasibility ($0\le u_t\le \alpha x_t$), and state-dependent
        cost functions, and we introduce \emph{Simplex Disturbance-Action Control}
        (SDAC), a DAC-inspired simplex-constrained policy class to compete with. The
        SDAC \emph{enforces feasibility} by construction while preserving DAC-style
        expressiveness.
    \item \textbf{Algorithmic and analytical contribution.} We develop an
        OMD controller with feasibility projection over SDAC parameters that
        \emph{enforces} the state-dependent constraint on the
        online control action, and we analyze regret against the best fixed SDAC
        policy in hindsight, proving a bound that is sublinear in the horizon and
        logarithmic in the policy dimension.
\end{itemize}

\subsection{Paper Organization}
Section~\ref{sec:model-setup} formalizes the adversarial storage-control model,
assumptions, and regret benchmark. Section~\ref{sec:sdac} introduces the SDAC policy
class and summarizes its structural properties. Section~\ref{sec:online-algorithm}
presents the online controller, Section~\ref{sec:regret-analysis} develops regret
guarantees, Section~\ref{sec:extensions} sketches immediate extensions, and
Section~\ref{sec:conclusion} concludes. Deferred proofs appear in the
Appendix.


\section{Notation}
\label{sec:notation}

We write $[z]_+ \triangleq \max\{z,0\}$. For $v\in\mathbb{R}^n$, the $i$-th
coordinate is $v^{(i)}$, and $\mathbf{1}$ denotes the all-ones vector (dimension
clear from context). The Euclidean inner product is $\langle\cdot,\cdot\rangle$. Time
indices may be nonpositive, with corresponding quantities taken to be zero
(e.g.\ $w_\tau\equiv 0$ for $\tau\le 0$). We write $\mathcal{A}$ for the online
algorithm and $\pi$ (or $\pi(M)$) for a fixed policy.


\section{Model and Problem Setup}
\label{sec:model-setup}

We consider a discrete-time process with times $t\in\{1,\ldots,T\}$. The system
state $x_t\in\mathbb{R}_{\ge 0}$ denotes the resource level at the beginning of
time $t$. The initial state is given as
\begin{equation}
\label{eq:model-init}
x_1 = 0.
\end{equation}
Fix a state-retention factor $\alpha\in(0,1)$. At each time, the controller chooses a
control action $u_t\in\mathbb{R}_{\ge 0}$ subject to the state constraint:
\begin{equation}
\label{eq:model-feasibility}
0 \le u_t \le \alpha x_t.
\end{equation}
The system then receives an arrival $w_t\in[0,1]$ and evolves according to
\begin{equation}
\label{eq:model-dynamics}
x_{t+1} = \alpha x_t - u_t + w_t.
\end{equation}

After the control action is selected, a convex cost function
\[
c_t: \mathbb{R}_{\ge 0}^2 \to \mathbb{R}
\]
is revealed, and the incurred cost is $c_t(x_t,u_t)$. We assume each $c_t$ is convex
in $(x,u)$ and satisfies a Lipschitz condition on the relevant domain. There exist
constants $L_x,L_u>0$ such that for all admissible $(x,u)$ and $(x',u')$,
\begin{equation}
\label{eq:model-lipschitz}
|c_t(x,u)-c_t(x',u')| \le L_x|x-x'| + L_u|u-u'|.
\end{equation}

The information pattern is adversarial and non-anticipatory. At time $t$, the
controller chooses $u_t$ before observing $w_t$ and $c_t$. Hence both the arrival
sequence $\{w_t\}_{t=1}^T$ and the sequence of cost functions $\{c_t\}_{t=1}^T$ may be
arbitrary.

\paragraph{Assumptions}

\begin{enumerate}
    \item \textbf{Retention factor:} $\alpha\in(0,1)$ is fixed and known (a fraction
        $1-\alpha$ of the state is lost each period).
    \item \textbf{Arrival bounds:} $w_t\in[0,1]$ for all $t\in\{1,\ldots,T\}$.
        In particular, arrivals are nonnegative resource increments, in contrast to
        signed process-noise disturbances common in classical control.
    \item \textbf{Adversarial timing:} at each time $t\in\{1,\ldots,T\}$, $u_t$ is chosen before observing
        $w_t$ and $c_t$.
    \item \textbf{Cost regularity:} each $c_t$ is convex in $(x,u)$ and obeys
        \eqref{eq:model-lipschitz}.
    \item \textbf{Feasible policy class:} both the online policy to be developed and the benchmark
     policies described below are required to satisfy \eqref{eq:model-feasibility} under dynamics
        \eqref{eq:model-dynamics}.
\end{enumerate}
We focus on $\alpha\in(0,1)$; the case $\alpha\ge 1$ is discussed in
Section~\ref{sec:extensions}.

Let $\{(x_t^{\mathcal{A}},u_t^{\mathcal{A}})\}_{t=1}^T$ denote the trajectory of an
online algorithm $\mathcal{A}$. Define its cumulative cost by
\begin{equation}
\label{eq:def-cumulative-cost}
J_T(\mathcal{A}) \triangleq \sum_{t=1}^T c_t\bigl(x_t^{\mathcal{A}},u_t^{\mathcal{A}}\bigr).
\end{equation}
For a fixed positive integer $H$, let $\pi$ range over the policy class
$\mathrm{SDAC}_H$ (defined in Section~\ref{sec:sdac}). For each
$\pi\in\mathrm{SDAC}_H$, let $\{(x_t^{\pi},u_t^{\pi})\}_{t=1}^T$ be the induced
feasible trajectory under the same arrival sequence. Regret against the best
fixed SDAC policy is defined by
\begin{equation}
\label{eq:def-regret}
\RegSDAC{T}
\triangleq
\sum_{t=1}^T c_t\bigl(x_t^{\mathcal{A}},u_t^{\mathcal{A}}\bigr)
-
\min_{\pi\in\mathrm{SDAC}_H}
\sum_{t=1}^T c_t\bigl(x_t^{\pi},u_t^{\pi}\bigr).
\end{equation}
While the online controller may update its SDAC parameters over time, the benchmark
in \eqref{eq:def-regret} is the best single fixed SDAC policy chosen in hindsight.
The objective is to design an online policy that preserves feasibility in
\eqref{eq:model-feasibility} for all $t\in\{1,\ldots,T\}$ and achieves sublinear regret as defined in
\eqref{eq:def-regret}.

\section{Simplex Disturbance-Action Control (SDAC)}
\label{sec:sdac}

\subsection{From Disturbance-Action Control to SDAC}
Disturbance-Action Control (DAC) parameterizes the control action as a linear
function of a finite window of past arrivals. This representation is central in
adversarial online control because the cost is convex in term of the parameters,
yields robust regret guarantees, and enables efficient online convex-optimization
updates. In our setting, a representative DAC form is
\begin{equation}
\label{eq:dac-param}
u_t = \sum_{i=1}^{H} M^{(i)} w_{t-i},
\end{equation}
where $H\in\mathbb{N}$ is the memory horizon and $M\in\mathbb{R}^H$ is a linear
filter. We allow nonpositive indices in such sums, with the convention $w_\tau\equiv
0$ for $\tau\le 0$.

For our constrained model, unconstrained DAC is insufficient. In particular, the
control action can violate \eqref{eq:model-feasibility}, either by having $u_t>\alpha x_t$ or $u_t<0$. We
therefore introduce \emph{Simplex Disturbance-Action Control} (SDAC), which replaces
unconstrained DAC filters by a simplex-constrained family tailored to this dynamics.
SDAC is designed to keep the DAC advantages above (adversarial robustness,
regret-based comparison, and tractable online updates) while adding feasibility by
construction.


\subsection{SDAC Definition}
Fix memory horizon $H\in\mathbb{N}$. Let
\begin{equation}
\label{eq:def-simplex}
\Delta_H \triangleq \left\{M\in\mathbb{R}_{\ge 0}^H:\; \sum_{i=1}^H
M^{(i)}\le1\right\}
\end{equation}
be the sub-probability simplex. For any $M\in\Delta_H$, define the policy $\pi(M)$ by
\begin{equation}
\label{eq:def-sdac-policy}
u_t^{\pi(M)} \;\triangleq\; \sum_{i=1}^{H} M^{(i)}\alpha^i w_{t-i},
\qquad w_\tau\equiv 0\;\text{for}\;\tau\le 0.
\end{equation}
For each time $t\in\{1,2,\ldots,T\}$, let
\begin{equation}
\label{eq:def-w-tilde}
\widetilde{w}_t \triangleq \big(\alpha w_{t-1},\alpha^2 w_{t-2},\ldots,\alpha^H w_{t-H}\big)^\top \in \mathbb{R}^H,
\end{equation}
with the same convention $w_\tau\equiv 0$ for $\tau\le 0$, so that
\eqref{eq:def-sdac-policy} can be written as
\begin{equation}
\label{eq:def-sdac-policy-vector}
u_t^{\pi(M)} = \langle M,\widetilde{w}_t\rangle.
\end{equation}

The induced state trajectory satisfies
\begin{equation}
\label{eq:sdac-state-update}
x_{t+1}^{\pi(M)} = \alpha x_t^{\pi(M)} - u_t^{\pi(M)} + w_t,
\qquad x_1^{\pi(M)}=0.
\end{equation}
The weights $\alpha^i$ align the arrival filter with the retention factor
$\alpha$ in \eqref{eq:model-dynamics}.


We now establish the two structural properties that make SDAC useful: the induced
control action always satisfies the feasibility constraint
\eqref{eq:model-feasibility}, and the induced cost is convex and Lipschitz in $M$, so
that regret minimization against the best fixed SDAC policy is an instance of online
convex optimization over the simplex. Both properties follow from an explicit
decomposition of the state trajectory that isolates its dependence on $M$.

Recall that we set $w_\tau\equiv 0$ for $\tau\le 0$, which is used in the base-case
arguments below.

\begin{lemma}[State decomposition under SDAC]
\label{lem:appendix-state-decomp}
Fix $H\in\mathbb{N}$ and $M\in\Delta_H$, and write $W\triangleq\sum_{i=1}^H
M^{(i)}\le 1$. Let $\pi(M)$ be defined by \eqref{eq:def-sdac-policy} and
\eqref{eq:sdac-state-update}. Then for all $t\in\{2,\ldots,T\}$,
\begin{equation}
\label{eq:appendix-state-decomp}
x_t^{\pi(M)}
=
\sum_{i=1}^H M^{(i)}\psi_t^{(i)}
+\bigl(1-W\bigr)\rho_t
=
\langle M,\psi_t\rangle+(1-W)\rho_t,
\end{equation}
where
\[
\psi_t^{(i)} \triangleq \sum_{j=1}^{i}\alpha^{j-1}w_{t-j}
\quad(i=1,\ldots,H),
\]
and
\[
\rho_t \triangleq \sum_{j=1}^{t-1}\alpha^{j-1}w_{t-j}.
\]
Equivalently, $\rho_t$ is the state $x_t^{\pi(0)}$ produced by the null policy $M=0$
(release nothing, $x_{\tau+1}=\alpha x_\tau+w_\tau$), and the coefficients $\psi_t^{(i)}$ and $\rho_t$ satisfy
$0\le \psi_t^{(i)}\le \rho_t\le 1/(1-\alpha)$. In particular, the map $M\mapsto x_t^{\pi(M)}$ is affine.
\end{lemma}

\begin{proof}
Isolating the $j=1$ term in the definition
of $\rho_{t+1}$ and reindexing gives the one-step identity
$\rho_{t+1}=\alpha\rho_t+w_t$ (with $\rho_1=0$). We prove
\eqref{eq:appendix-state-decomp} by induction on $t$.

\emph{Base case $t=2$.} Since $x_1^{\pi(M)}=0$ and $w_{1-i}=0$ for all $i\ge 1$, the
policy \eqref{eq:def-sdac-policy} gives $u_1^{\pi(M)}=\sum_{i=1}^H M^{(i)}\alpha^i
w_{1-i}=0$, and the update \eqref{eq:sdac-state-update} gives $x_2^{\pi(M)}=\alpha
x_1^{\pi(M)}-u_1^{\pi(M)}+w_1=w_1$. On the other hand $\psi_2^{(i)}=w_1$ for every
$i$ and $\rho_2=w_1$, so $\langle M,\psi_2\rangle+(1-W)\rho_2=Ww_1+(1-W)w_1=w_1$,
matching \eqref{eq:appendix-state-decomp}.

\emph{Inductive step.} Assume \eqref{eq:appendix-state-decomp} holds at time
$t\in\{2,\ldots,T-1\}$. Using \eqref{eq:def-sdac-policy}, \eqref{eq:sdac-state-update}, and the induction
hypothesis,
\[
\begin{aligned}
x_{t+1}^{\pi(M)}
&= \alpha x_t^{\pi(M)} - \sum_{i=1}^H M^{(i)}\alpha^i w_{t-i} + w_t \\
&= \Bigl(\alpha\langle M,\psi_t\rangle - \sum_{i=1}^H M^{(i)}\alpha^i w_{t-i}\Bigr) \\
&\qquad + \alpha(1-W)\rho_t + w_t .
\end{aligned}
\]
Reindexing the inner summation ($j\mapsto j+1$) yields
\[
\begin{aligned}
&\alpha\langle M,\psi_t\rangle-\sum_{i=1}^H M^{(i)}\alpha^i w_{t-i} \\
&\qquad=\sum_{i=1}^H M^{(i)}\!\left(\sum_{j=1}^{i-1}\alpha^j w_{t-j}\right) \\
&\qquad=\sum_{i=1}^H M^{(i)}\bigl(\psi_{t+1}^{(i)}-w_t\bigr)
=\langle M,\psi_{t+1}\rangle-W w_t ,
\end{aligned}
\]
where the second equality uses $\psi_{t+1}^{(i)}=w_t+\sum_{j=1}^{i-1}\alpha^j
w_{t-j}$. Substituting this together with $\rho_{t+1}=\alpha\rho_t+w_t$,
\[
\begin{aligned}
x_{t+1}^{\pi(M)}
&=\langle M,\psi_{t+1}\rangle - W w_t + \alpha(1-W)\rho_t + w_t \\
&=\langle M,\psi_{t+1}\rangle + (1-W)\bigl(w_t+\alpha\rho_t\bigr) \\
&=\langle M,\psi_{t+1}\rangle+(1-W)\rho_{t+1},
\end{aligned}
\]
which is \eqref{eq:appendix-state-decomp} at time $t+1$.

\emph{Bounds.} Each $\psi_t^{(i)}$ is the truncation to $j\le i$ of the nonnegative
series defining $\rho_t$, so $0\le\psi_t^{(i)}\le\rho_t$; and since $w_\tau\in[0,1]$,
the geometric series gives
$\rho_t=\sum_{j=1}^{t-1}\alpha^{j-1}w_{t-j}\le\sum_{j=1}^{\infty}\alpha^{j-1}=1/(1-\alpha)$.
Affineness in $M$ is immediate, as $\psi_t$ and $\rho_t$ do not depend on $M$.
\end{proof}

The first consequence of this decomposition is that feasibility is preserved within
the class.

\begin{lemma}[Feasibility of SDAC policies]
\label{lem:appendix-feasibility}
For every fixed $M\in\Delta_H$, the trajectory induced by
\eqref{eq:def-sdac-policy}--\eqref{eq:sdac-state-update} satisfies
\[
0\le u_t^{\pi(M)}\le \alpha x_t^{\pi(M)} \quad \text{for all } t\in\{1,\ldots,T\}.
\]
\end{lemma}

\begin{proof}
Nonnegativity of $u_t^{\pi(M)}$ follows from \eqref{eq:def-sdac-policy}, because
$M^{(i)}\ge 0$, $\alpha^i\ge 0$, and $w_{t-i}\ge 0$ for every $i$.

For the upper bound at $t=1$, we have $x_1^{\pi(M)}=0$ and $w_{1-i}=0$ for all $i\ge
1$, so $u_1^{\pi(M)}=0=\alpha x_1^{\pi(M)}$.

Now fix $t\in\{2,\ldots,T\}$. Lemma~\ref{lem:appendix-state-decomp} gives $x_t^{\pi(M)}=\langle
M,\psi_t\rangle+(1-W)\rho_t$ with $W=\sum_{i=1}^H M^{(i)}\le 1$. Since
$(1-W)\rho_t\ge 0$,
\[
\begin{aligned}
\alpha x_t^{\pi(M)}
&\ge \alpha\langle M,\psi_t\rangle
= \alpha\sum_{i=1}^H M^{(i)}\!\left(\sum_{j=1}^{i}\alpha^{j-1}w_{t-j}\right) \\
&= \sum_{i=1}^H M^{(i)}\!\left(\sum_{j=1}^{i}\alpha^j w_{t-j}\right).
\end{aligned}
\]
Collecting the coefficient of each $w_{t-i}$ and using that all parameters are
nonnegative, for every $i$ we have $\sum_{j=i}^{H} M^{(j)} \ge M^{(i)}$, so
\[
\begin{aligned}
\alpha x_t^{\pi(M)}
&\ge
\sum_{i=1}^{H}\left(\sum_{j=i}^{H}M^{(j)}\right)\alpha^i w_{t-i} \\
&\ge
\sum_{i=1}^{H}M^{(i)}\alpha^i w_{t-i}
=
u_t^{\pi(M)},
\end{aligned}
\]
since $w_{t-i}\ge 0$.
\end{proof}

The second consequence is that SDAC is compatible with online convex optimization:
for fixed arrival sequences, $u_t^{\pi(M)}$ is linear in $M$ by
\eqref{eq:def-sdac-policy-vector}, and $x_t^{\pi(M)}$ is affine in $M$ by
Lemma~\ref{lem:appendix-state-decomp}. Since each $c_t$ is convex in $(x,u)$, this
already shows the induced cost is convex in $M$; the next lemma also quantifies its
Lipschitz constant, which we use later to tune the online algorithm's step size.

\begin{lemma}[Convexity and Lipschitz continuity of surrogate cost]
\label{lem:appendix-surrogate}
Define
\[
\ell_t(M) \triangleq c_t\bigl(x_t^{\pi(M)},u_t^{\pi(M)}\bigr), \qquad M\in\Delta_H.
\]
Then $\ell_t$ is convex on $\Delta_H$, and for all $M,M'\in\Delta_H$,
\[
|\ell_t(M)-\ell_t(M')|
\le
\left(\frac{L_x}{1-\alpha}+\alpha L_u\right)\|M-M'\|_1.
\]
\end{lemma}

\begin{proof}
By Lemma~\ref{lem:appendix-state-decomp} and
\eqref{eq:def-sdac-policy-vector}, with $W=\sum_{i=1}^H M^{(i)}$,
\[
\ell_t(M)=c_t\bigl(\langle M,\psi_t\rangle+(1-W)\rho_t,\ \langle
M,\widetilde{w}_t\rangle\bigr).
\]
The map $M\mapsto\bigl(x_t^{\pi(M)},u_t^{\pi(M)}\bigr)$ is affine
(Lemma~\ref{lem:appendix-state-decomp} and
\eqref{eq:def-sdac-policy-vector}) and $c_t$ is convex, so $\ell_t$ is convex.

For the Lipschitz estimate, the decomposition gives, for all $M,M'\in\Delta_H$,
\[
\begin{aligned}
x_t^{\pi(M)}-x_t^{\pi(M')}
&=\langle M-M',\psi_t\rangle-(W-W')\rho_t \\
&=\langle M-M',\,\psi_t-\rho_t\mathbf{1}\rangle,
\end{aligned}
\]
using $W-W'=\langle M-M',\mathbf{1}\rangle$, while
$u_t^{\pi(M)}-u_t^{\pi(M')}=\langle M-M',\widetilde{w}_t\rangle$. Hence, by
\eqref{eq:model-lipschitz} and Hölder's inequality,
\begin{align*}
|\ell_t(M)-\ell_t(M')|
&\le L_x\bigl|\langle M-M',\,\psi_t-\rho_t\mathbf{1}\rangle\bigr| \\
&\quad + L_u\bigl|\langle M-M',\widetilde{w}_t\rangle\bigr| \\
&\le \bigl(L_x\,\|\psi_t-\rho_t\mathbf{1}\|_\infty + L_u\,\|\widetilde{w}_t\|_\infty\bigr) \\
&\qquad\cdot\|M-M'\|_1.
\end{align*}
Since $0\le\psi_t^{(i)}\le\rho_t\le 1/(1-\alpha)$
(Lemma~\ref{lem:appendix-state-decomp}), we have
$\|\psi_t-\rho_t\mathbf{1}\|_\infty=\max_i\bigl(\rho_t-\psi_t^{(i)}\bigr)\le\rho_t\le
1/(1-\alpha)$, and $\|\widetilde{w}_t\|_\infty=\max_i\alpha^i w_{t-i}\le\alpha$
because $w_\tau\in[0,1]$ and $i\ge 1$. This yields the claim, and one may take the
surrogate Lipschitz constant
\[
L := \alpha L_u + \frac{L_x}{1-\alpha},
\]
as used in the main text.
\end{proof}

Regret minimization against the best fixed SDAC policy thus becomes online convex
optimization over the simplex.

\subsection{Expressiveness: Linear State-Feedback Approximation}
\label{subsec:sdac-expressiveness}

A standard policy class in online control is the family of linear state-feedback
policies $u_t=-Kx_t$, and a foundational property of Disturbance-Action Control is
that the DAC class \emph{approximates} any strongly stable linear policy, with
approximation error decaying geometrically in the memory horizon $H$
\cite{agarwal2019online,hazan_introduction_2025}.

We establish the analogue of this guarantee for SDAC. In our scalar,
state-constrained model the feasible linear state-feedback policies are exactly the
proportional (constant-fraction) control actions $u_t=q\,x_t$ with gain
$q\in[0,\alpha]$: since $x_t\ge 0$, the feasibility constraint $0\le u_t\le\alpha
x_t$ in \eqref{eq:model-feasibility} holds if and only if $q\in[0,\alpha]$. The next
lemma shows that every such policy with $q\in(0,\alpha)$ is approximated,
\emph{within the feasible class} $\mathrm{SDAC}_H$, by explicit parameters, with
cost approximation error that decays geometrically in $H$.

\begin{lemma}[SDAC approximates linear state-feedback policies]
\label{lem:sdac-linear-approx}
Fix $q\in(0,\alpha)$, set $\beta\triangleq\alpha-q$, and let $\pi_q^{\mathrm{lin}}$
denote the feasible linear policy
\[
u_t^{\mathrm{lin}}=q\,x_t^{\mathrm{lin}},
\qquad
x_{t+1}^{\mathrm{lin}}=\beta\,x_t^{\mathrm{lin}}+w_t,
\qquad
t\in\{1,\ldots,T\},
\qquad
x_1^{\mathrm{lin}}=0.
\]
Define $M_q^\star\in\mathbb{R}^H$ by
\[
M_q^{\star(i)}\triangleq\frac{q}{\alpha}\Bigl(\frac{\beta}{\alpha}\Bigr)^{i-1},
\qquad i=1,\ldots,H.
\]
Then $M_q^\star\in\Delta_H$, and if each $c_t$ satisfies
\eqref{eq:model-lipschitz}, then for every $\{w_t\}\subset[0,1]$,
\begin{equation}
\label{eq:sdac-lin-approx-cost}
\begin{aligned}
&\Bigl|\sum_{t=1}^T c_t\bigl(x_t^{\pi(M_q^\star)},u_t^{\pi(M_q^\star)}\bigr)
-\sum_{t=1}^T c_t\bigl(x_t^{\mathrm{lin}},u_t^{\mathrm{lin}}\bigr)\Bigr| \\
&\qquad\le
\Bigl(L_u+\frac{L_x}{1-\alpha}\Bigr)\frac{q\,\beta^H}{1-\beta}\,T.
\end{aligned}
\end{equation}
In particular, the cost error decays geometrically in $H$, so $\mathrm{SDAC}_H$
approximates $\pi_q^{\mathrm{lin}}$ arbitrarily well as $H\to\infty$. The proof is
deferred to Appendix~\ref{sec:appendix-proofs}.
\end{lemma}

This expressiveness property has a direct algorithmic consequence. Because the regret
bound of Theorem~\ref{thm:regret-bound} depends on the memory horizon only through
$\sqrt{\log(H+1)}$, the window $H$ can be enlarged cheaply; combining
Lemma~\ref{lem:sdac-linear-approx} with that bound then shows the online controller
competes with the best fixed linear policy. We make this precise in
Corollary~\ref{cor:regret-vs-linear}, mirroring the role of DAC-to-linear
approximation in online control \cite{agarwal2019online}.



 

\section{Online Algorithm}
\label{sec:online-algorithm}

We now present our online algorithm, which updates SDAC parameters by
online mirror descent (OMD) with unnormalized negentropy (equivalently, exponentiated gradient)
over the sub-probability simplex $\Delta_H\subset\mathbb{R}^H$ with a fixed step size
$\eta>0$.

For each time $t\in\{1,\ldots,T\}$, let $M_t\in\Delta_H$ be the online parameter. We initialize
at
\begin{equation}
\label{eq:alg-init}
M_1^{(i)}\triangleq\frac{1}{H+1},\qquad i=1,\ldots,H
\end{equation}
so that the Bregman diameter of $\Delta_H$ relative to $M_1$ under unnormalized
negentropy is at most $\log(H+1)$ (see the proof of Theorem~\ref{thm:regret-bound}).

\paragraph{Control-action selection}
The controller first forms the unconstrained SDAC control action
\begin{equation}
\label{eq:alg-unconstrained-action}
\widehat{u}_t \triangleq \langle M_t,\widetilde{w}_t\rangle,
\end{equation}
then applies a feasibility projection onto the admissible action interval $[0,\alpha
x_t^{\mathcal{A}}]$:
\begin{equation}
\label{eq:alg-feasible-action}
u_t^{\mathcal{A}} \triangleq \min\{\widehat{u}_t,\,\alpha x_t^{\mathcal{A}}\}.
\end{equation}
The state update is
\begin{equation}
\label{eq:alg-state-update}
x_{t+1}^{\mathcal{A}} = \alpha x_t^{\mathcal{A}} - u_t^{\mathcal{A}} + w_t,
\qquad x_1^{\mathcal{A}}=0.
\end{equation}
By construction, $\widehat{u}_t\ge 0$ and $u_t^{\mathcal{A}}\in[0,\alpha x_t^{\mathcal{A}}]$;
with $x_1^{\mathcal{A}}=0$ and $w_t\in[0,1]$, induction on
\eqref{eq:alg-state-update} yields $x_t^{\mathcal{A}}\ge 0$ for all $t\in\{1,\ldots,T\}$,
so the online trajectory satisfies \eqref{eq:model-feasibility}.

\paragraph{Surrogate cost and OMD update}
Define the surrogate cost at each time $t\in\{1,\ldots,T\}$ by
\begin{equation}
\label{eq:alg-surrogate-cost}
\ell_t(M) \triangleq c_t\bigl(x_t^{\pi(M)},u_t^{\pi(M)}\bigr), \qquad M\in\Delta_H,
\end{equation}
which is convex and Lipschitz by Lemma~\ref{lem:appendix-surrogate}. Let
$g_t\in\partial \ell_t(M_t)$ be a subgradient. Let
\begin{equation}
\label{eq:def-negentropy}
\Psi(M) \triangleq \sum_{i=1}^H\bigl(M^{(i)}\log M^{(i)}-M^{(i)}\bigr), \qquad
M\in\mathbb{R}_{\ge0}^H,
\end{equation}
denote the unnormalized negentropy (with the convention $0\log0\triangleq0$), and let
\[
D_\Psi(M\|V)\triangleq\sum_{i=1}^H\Bigl(M^{(i)}\log\tfrac{M^{(i)}}{V^{(i)}}-M^{(i)}+V^{(i)}\Bigr)
\]
be its Bregman divergence, the generalized Kullback--Leibler divergence.

We define the OMD update
\begin{equation}
\label{eq:alg-omd-argmin}
M_{t+1} \triangleq \operatorname*{arg\,min}_{M\in\Delta_H}
\Bigl\{\eta\langle g_t,M\rangle + D_\Psi(M\|M_t)\Bigr\}.
\end{equation}
All of the regret analysis in Section~\ref{sec:regret-analysis} works directly with
this optimization-based definition.

\paragraph{Closed-form solution}
Although \eqref{eq:alg-omd-argmin} is a constrained optimization problem, it is easy
to solve analytically because $\Psi$ is separable and the objective is linear in
$g_t$. As shown in Lemma~\ref{lem:omd-closed-form}, the minimizer in
\eqref{eq:alg-omd-argmin} is exactly the result of two elementary steps: a
multiplicative step,
\begin{equation}
\label{eq:alg-omd-mult}
\widetilde{M}_{t+1}^{(i)} \triangleq M_t^{(i)}\exp\!\bigl(-\eta g_t^{(i)}\bigr), \qquad i=1,\ldots,H,
\end{equation}
followed by the Bregman projection onto $\Delta_H$, triggered only when the resulting
mass exceeds one:
\begin{equation}
\label{eq:alg-omd-proj}
M_{t+1}^{(i)}
\triangleq
\frac{\widetilde{M}_{t+1}^{(i)}}{\max\!\bigl(1,S_{t+1}\bigr)},
\quad
S_{t+1}\triangleq\sum_{j=1}^H \widetilde{M}_{t+1}^{(j)},
\end{equation}
for $i=1,\ldots,H$. This is the form we implement in practice.

This yields a feasible online trajectory satisfying \eqref{eq:model-feasibility} by
\eqref{eq:alg-feasible-action} and performs OMD over $\Delta_H$
through \eqref{eq:alg-surrogate-cost} and \eqref{eq:alg-omd-argmin}.


\section{Regret Analysis}
\label{sec:regret-analysis}
The following theorem states that the online controller achieves sublinear regret
 against every fixed SDAC policy $\pi(M)$, where $M$ is any parameter in $\Delta_H$.
The regret proof has two parts: a standard online convex optimization bound for the
surrogate costs \(\ell_t\), and a stability bound showing that the
feasibility-projected online control action stays close to the SDAC policy trajectory
driven by the same parameters. Some technical lemmas required for the proof are deferred to
Appendix~\ref{sec:appendix-proofs}.

\begin{theorem}[Regret bound]
\label{thm:regret-bound}
Suppose each cost function \(c_t\) is convex and satisfies
\eqref{eq:model-lipschitz}. Then the OMD controller with
feasibility projection from Section~\ref{sec:online-algorithm}, initialized as in
\eqref{eq:alg-init} and using parameter step size $\eta>0$, satisfies
\[
\begin{aligned}
\RegSDAC{T}
&\le
\frac{\log(H+1)}{\eta}
+
\eta L^2 T\Bigl(\tfrac{1}{2}+\tfrac{\alpha}{(1-\alpha)^2}\Bigr), \\
L &:= \alpha L_u + \tfrac{L_x}{1-\alpha}.
\end{aligned}
\]
In particular, the choice
\[
\eta^*
=
\sqrt{\frac{\log(H+1)}{L^2 T\bigl(\tfrac{1}{2}+\tfrac{\alpha}{(1-\alpha)^2}\bigr)}}
\]
gives
\[
\RegSDAC{T}
\le
2L\sqrt{\Bigl(\tfrac{1}{2}+\tfrac{\alpha}{(1-\alpha)^2}\Bigr)T\log(H+1)}.
\]
\end{theorem}


Combining Theorem~\ref{thm:regret-bound} with the approximation guarantee of
Lemma~\ref{lem:sdac-linear-approx} shows that the controller is competitive not only
with the best fixed SDAC policy but with the entire family of feasible linear
state-feedback policies.

\begin{corollary}[No-regret against linear state-feedback policies]
\label{cor:regret-vs-linear}
Run the OMD controller with feasibility projection of
Section~\ref{sec:online-algorithm} with the step size $\eta^*$ of
Theorem~\ref{thm:regret-bound}. Fix any gain $q\in(0,\alpha)$, set
$\beta\triangleq\alpha-q$, and let
$\bigl(x_t^{\mathrm{lin}},u_t^{\mathrm{lin}}\bigr)$ be the trajectory of the linear
policy $\pi_q^{\mathrm{lin}}$ from Lemma~\ref{lem:sdac-linear-approx}. Define
\[
\RegLin{T}
\triangleq
\sum_{t=1}^T c_t\bigl(x_t^{\mathcal{A}},u_t^{\mathcal{A}}\bigr)
-\sum_{t=1}^T c_t\bigl(x_t^{\mathrm{lin}},u_t^{\mathrm{lin}}\bigr).
\]
Then
\[
\begin{aligned}
\RegLin{T}
&\le
2L\sqrt{\Bigl(\tfrac12+\tfrac{\alpha}{(1-\alpha)^2}\Bigr)T\log(H+1)} \\
&\quad +\Bigl(L_u+\frac{L_x}{1-\alpha}\Bigr)\frac{q\,\beta^H}{1-\beta}\,T,
\end{aligned}
\]
where $L=\alpha L_u+\tfrac{L_x}{1-\alpha}$. In particular, choosing the memory
horizon $H=\bigl\lceil \log T/\log(1/\beta)\bigr\rceil$ forces $\beta^H\le 1/T$, so
the second term is $O(1)$ and
\[
\RegLin{T}
=
O\!\left(\sqrt{T\log\log T}\,\right).
\]
Thus $\RegLin{T}=o(T)$ against every fixed linear state-feedback policy
$u_t=q\,x_t$, $q\in(0,\alpha)$: the controller is no-regret relative to this class.
\end{corollary}

\begin{proof}
By the regret definition \eqref{eq:def-regret} and the optimized bound of
Theorem~\ref{thm:regret-bound},
\[
\begin{aligned}
\sum_{t=1}^T c_t\bigl(x_t^{\mathcal{A}},u_t^{\mathcal{A}}\bigr)
&\le
\min_{M\in\Delta_H}\sum_{t=1}^T c_t\bigl(x_t^{\pi(M)},u_t^{\pi(M)}\bigr) \\
&\quad +2L\sqrt{\Bigl(\tfrac12+\tfrac{\alpha}{(1-\alpha)^2}\Bigr)T\log(H+1)}.
\end{aligned}
\]
Since $M_q^\star\in\Delta_H$ by Lemma~\ref{lem:sdac-linear-approx}, the minimum is no
larger than the cost of $\pi(M_q^\star)$, which by \eqref{eq:sdac-lin-approx-cost}
obeys
\[
\begin{aligned}
&\sum_{t=1}^T c_t\bigl(x_t^{\pi(M_q^\star)},u_t^{\pi(M_q^\star)}\bigr) \\
&\qquad\le
\sum_{t=1}^T c_t\bigl(x_t^{\mathrm{lin}},u_t^{\mathrm{lin}}\bigr)
+\Bigl(L_u+\frac{L_x}{1-\alpha}\Bigr)\frac{q\,\beta^H}{1-\beta}\,T.
\end{aligned}
\]
Chaining the two displays gives, by the definition of $\RegLin{T}$ above, the stated
bound. For the asymptotic statement,
$H=\bigl\lceil \log T/\log(1/\beta)\bigr\rceil$ gives $\beta^H\le 1/T$, so the second
term is at most $\bigl(L_u+\tfrac{L_x}{1-\alpha}\bigr)\tfrac{q}{1-\beta}=O(1)$, while
$\log(H+1)=O(\log\log T)$ in the first term.
\end{proof}


\section{Extensions}
\label{sec:extensions}

We briefly sketch three directions in which the present framework extends with only
minor changes to the policy class and algorithm. A full analysis is left for future
work.

\paragraph{Retention $\alpha\ge 1$.}
Extending SDAC and our algorithm and analysis to handle $\alpha\ge 1$ is
straightforward: one can augment the SDAC action by a fixed proportional drain,
replacing $u_t^{\pi(M)}=\langle M,\widetilde{w}_t\rangle$ with
$u_t^{\pi(M)}=Kx_t+\langle M,\widetilde{w}_t\rangle$ for a fixed $K$ satisfying
$0<\alpha-K<1$, which reduces the dynamics to an equivalent system with retention
in $(0,1)$. We deliberately restricted attention to the simpler $\alpha\in(0,1)$
environment for elegance and to avoid extra lines of derivation.

\paragraph{Vector action / multi-task allocation.}
Consider a single storage state $x_t$ with a vector action
$u_t\in\mathbb{R}_{\ge 0}^n$ that allocates resource across $n$ tasks, subject to
$\mathbf{1}^\top u_t\le\alpha x_t$, so that the total drain from storage is the sum of
the task allocations. The dynamics remain
$x_{t+1}=\alpha x_t-\mathbf{1}^\top u_t+w_t$, and the cost $c_t(x_t,u_t)$ is convex in
$(x,u)$. SDAC lifts by replacing the filter vector $M\in\Delta_H$ with a matrix
$M\in\mathbb{R}_{\ge 0}^{n\times H}$ whose entries sum to at most one (a matrix
simplex), with each row corresponding to a task and each column to a memory lag.
OMD then runs over this matrix simplex with the same feasibility-projection idea as
in the scalar case.

\paragraph{Multiple storages with diagonal dynamics.}
Suppose the state is $x_t\in\mathbb{R}_{\ge 0}^d$ and evolves as
$x_{t+1}=A x_t-u_t+w_t$ with $A=\mathrm{diag}(\alpha_1,\ldots,\alpha_d)$,
$\alpha_i\in(0,1)$, arrivals $w_t\in[0,1]^d$, and per-coordinate constraints
$0\le u_t^{(i)}\le\alpha_i x_t^{(i)}$. Because $A$ is diagonal, the coordinates
decouple into $d$ independent scalar storage problems. The natural SDAC class is a
product of simplices, and OMD may be run independently on each coordinate. Cross-coupled
storages (non-diagonal $A$) fall outside this easy reduction.

These directions preserve feasibility by construction and the OMD structure developed
above; fleshing out the corresponding regret bounds is left for future work.

\section{Conclusion}
\label{sec:conclusion}

We studied adversarial online control of a storage system under  state-dependent
action  constraints. The SDAC policy class enforces
feasibility by construction, and the OMD controller with
feasibility projection over
SDAC parameters achieves sublinear regret $\RegSDAC{T}$ against the best fixed SDAC
policy, with only logarithmic dependence on the memory horizon. The same guarantees
extend to competition with feasible linear state-feedback policies, i.e. sublinear
$\RegLin{T}$, via the SDAC approximation result. The framework is directly applicable to energy-harvesting battery management
and related storage settings in which decisions must respect
state constraints in real time. Immediate lifts---retention $\alpha\ge 1$,
multi-task vector allocation, and diagonal multi-storage---are outlined in
Section~\ref{sec:extensions}.



\appendices
\section{Deferred Proofs}
\label{sec:appendix-proofs}

This appendix collects the deferred proof of the linear-state-feedback
approximation guarantee from Section~\ref{sec:sdac}, together with the OMD and
stability lemmas leading up to the regret bound stated in
Section~\ref{sec:regret-analysis}.

Recall that we set $w_\tau\equiv 0$ for $\tau\le 0$, which is used in several
base-case arguments below.

We first prove the linear-policy approximation guarantee stated in
Section~\ref{subsec:sdac-expressiveness}.

\begin{proof}[Proof of Lemma~\ref{lem:sdac-linear-approx}]
Fix $q\in(0,\alpha)$, set $\beta=\alpha-q\in(0,\alpha)$, and recall $w_\tau\equiv 0$
for $\tau\le 0$.

\emph{Feasibility and closed form of the linear policy.} Substituting
$u_t^{\mathrm{lin}}=q\,x_t^{\mathrm{lin}}$ into its state update gives the closed
loop 
\begin{equation}
\label{eq:model-closed-loop}
x_{t+1}^{\mathrm{lin}}=\alpha x_t^{\mathrm{lin}}-q x_t^{\mathrm{lin}}+w_t=\beta x_t^{\mathrm{lin}}+w_t,
\end{equation}
with $x_1^{\mathrm{lin}}=0$. Since $\beta>0$ and $w_t\ge 0$,
induction yields $x_t^{\mathrm{lin}}\ge 0$, and then $0\le u_t^{\mathrm{lin}}=q
x_t^{\mathrm{lin}}\le\alpha x_t^{\mathrm{lin}}$ because $q\le\alpha$; thus
$\pi_q^{\mathrm{lin}}$ satisfies \eqref{eq:model-feasibility}. Unrolling the
closed loop recursion \eqref{eq:model-closed-loop} from the initialization
$x_1^{\mathrm{lin}}=0$ (by \eqref{eq:model-init}), yields for all $t\in\{1,\ldots,T\}$,
\[
x_t^{\mathrm{lin}}=\sum_{j=1}^{t-1}\beta^{j-1}w_{t-j},
\quad \text{and} \quad
u_t^{\mathrm{lin}}=q\sum_{j=1}^{t-1}\beta^{j-1}w_{t-j}.
\]

\emph{Part 1.} Each $M_q^{\star(i)}=\tfrac{q}{\alpha}(\beta/\alpha)^{i-1}\ge 0$.
Using $1-\beta/\alpha=(\alpha-\beta)/\alpha=q/\alpha$ and the finite geometric sum,
\[
\begin{aligned}
\sum_{i=1}^H M_q^{\star(i)}
&=\frac{q}{\alpha}\cdot\frac{1-(\beta/\alpha)^H}{1-\beta/\alpha} \\
&=\frac{q}{\alpha}\cdot\frac{\alpha}{q}\bigl(1-(\beta/\alpha)^H\bigr)
=1-(\beta/\alpha)^H<1,
\end{aligned}
\]
since $0<\beta/\alpha<1$. Hence $M_q^\star\in\Delta_H$ and
$\pi(M_q^\star)\in\mathrm{SDAC}_H$.

\emph{Part 2.} By the SDAC definition \eqref{eq:def-sdac-policy} and the identity
$M_q^{\star(i)}\alpha^i=q\beta^{i-1}$,
\[
u_t^{\pi(M_q^\star)}=\sum_{i=1}^H M_q^{\star(i)}\alpha^i w_{t-i}=q\sum_{i=1}^H\beta^{i-1}w_{t-i}.
\]
Truncating with $w_{t-i}=0$ for $i\ge t$ yields
\[
u_t^{\pi(M_q^\star)}=q\sum_{i=1}^{\min\{H,t-1\}}\beta^{i-1}w_{t-i}.
\]
Subtracting from $u_t^{\mathrm{lin}}$ cancels the first $\min\{H,t-1\}$ terms:
\[
u_t^{\mathrm{lin}}-u_t^{\pi(M_q^\star)}=q\sum_{j=H+1}^{t-1}\beta^{j-1}w_{t-j}\ge 0,
\]
which equals $0$ when $t\le H+1$. Using $w_{t-j}\in[0,1]$ and summing the geometric
tail,
\[
u_t^{\mathrm{lin}}-u_t^{\pi(M_q^\star)}\le q\sum_{j=H+1}^{\infty}\beta^{j-1}=\frac{q\,\beta^H}{1-\beta}.
\]

\emph{Part 3.} Let $\delta_t\triangleq x_t^{\pi(M_q^\star)}-x_t^{\mathrm{lin}}$ and
$e_t\triangleq
u_t^{\mathrm{lin}}-u_t^{\pi(M_q^\star)}\in\bigl[0,\,q\beta^H/(1-\beta)\bigr]$ by Part
2. The two trajectories follow the same dynamics \eqref{eq:model-dynamics} driven by
the same arrival, so the arrival cancels on subtraction:
\[
\delta_{t+1}=\alpha\,\delta_t-\bigl(u_t^{\pi(M_q^\star)}-u_t^{\mathrm{lin}}\bigr)=\alpha\,\delta_t+e_t,
\qquad \delta_1=0.
\]
Unrolling gives $\delta_t=\sum_{s=1}^{t-1}\alpha^{t-1-s}e_s\ge 0$, and since $e_s\le
q\beta^H/(1-\beta)$ and $\sum_{s=1}^{t-1}\alpha^{t-1-s}\le 1/(1-\alpha)$,
\[
0\le\delta_t\le\frac{q\,\beta^H}{1-\beta}\cdot\frac{1}{1-\alpha}=\frac{q\,\beta^H}{(1-\beta)(1-\alpha)}.
\]

\emph{Part 4.} By the Lipschitz condition \eqref{eq:model-lipschitz} and Parts 2--3,
for each $t\in\{1,\ldots,T\}$,
\[
\begin{aligned}
&\bigl|c_t\bigl(x_t^{\pi(M_q^\star)},u_t^{\pi(M_q^\star)}\bigr)
-c_t\bigl(x_t^{\mathrm{lin}},u_t^{\mathrm{lin}}\bigr)\bigr| \\
&\qquad\le L_x\,|\delta_t|+L_u\,|e_t|
\le\Bigl(\frac{L_x}{1-\alpha}+L_u\Bigr)\frac{q\,\beta^H}{1-\beta}.
\end{aligned}
\]
Summing over $t=1,\ldots,T$ and applying the triangle inequality yields Part 4.
\end{proof}

The remaining lemmas establish the OMD and stability estimates used in the proof of
Theorem~\ref{thm:regret-bound}; they rely on the state-decomposition, feasibility,
and surrogate-cost lemmas already established in Section~\ref{sec:sdac}
(Lemmas~\ref{lem:appendix-state-decomp}, \ref{lem:appendix-feasibility}, and
\ref{lem:appendix-surrogate}).

\begin{lemma}[Bregman projection onto $\Delta_H$]
\label{lem:bregman-projection}
Let $\Psi$ and $D_\Psi$ be the unnormalized negentropy and its Bregman divergence
from \eqref{eq:def-negentropy}. Fix $\widetilde{M}\in\mathbb{R}_{>0}^H$ and set
$S\triangleq\sum_{i=1}^H \widetilde{M}^{(i)}$. Then the Bregman projection of
$\widetilde{M}$ onto $\Delta_H=\{M\in\mathbb{R}_{\ge0}^H:\sum_{i} M^{(i)}\le1\}$,
\[
P(\widetilde{M})\triangleq\operatorname*{arg\,min}_{M\in\Delta_H} D_\Psi(M\|\widetilde{M}),
\]
is given coordinatewise by
\[
P(\widetilde{M})^{(i)}=
\frac{\widetilde{M}^{(i)}}{\max(1,S)}.
\]
\end{lemma}

\begin{proof}
Since $\widetilde{M}$ is fixed, the terms $+\widetilde{M}^{(i)}$ are constant, so
$P(\widetilde{M})$ minimizes
\[
F(M)\triangleq\sum_{i=1}^H\Bigl(M^{(i)}\log\tfrac{M^{(i)}}{\widetilde{M}^{(i)}}-M^{(i)}\Bigr)
\]
over $\Delta_H$. The function $F$ is strictly convex on $\mathbb{R}_{\ge0}^H$ and
$\Delta_H$ is closed and convex, so the minimizer is unique; its partial derivatives
are $\partial F/\partial M^{(i)}=\log\bigl(M^{(i)}/\widetilde{M}^{(i)}\bigr)$.

\emph{Case $S\le 1$.} Here $\widetilde{M}\in\Delta_H$. Setting $\partial F/\partial
M^{(i)}=0$ gives the unconstrained minimizer $M=\widetilde{M}$ over
$\mathbb{R}_{\ge0}^H$, which is feasible; hence $P(\widetilde{M})=\widetilde{M}$.

\emph{Case $S>1$.} Here $\widetilde{M}\notin\Delta_H$, so the constraint $\sum_i
M^{(i)}\le1$ must be active at the optimum. As $\widetilde{M}^{(i)}>0$, the optimal
coordinates are strictly positive and the nonnegativity constraints are inactive;
introducing a multiplier $\lambda$ for $\sum_i M^{(i)}=1$, KKT stationarity reads
\[
\log\bigl(M^{(i)}/\widetilde{M}^{(i)}\bigr)+\lambda=0
\quad\Longrightarrow\quad
M^{(i)}=\widetilde{M}^{(i)}e^{-\lambda}.
\]
The active constraint $\sum_i M^{(i)}=1$ forces $e^{-\lambda}=1/S$, so
$M^{(i)}=\widetilde{M}^{(i)}/S$; moreover $\lambda=\log S>0$, so dual feasibility
holds. Equivalently, $P(\widetilde{M})^{(i)}=\widetilde{M}^{(i)}/\max(1,S)$ in both
cases.
\end{proof}

\begin{lemma}[Closed-form solution of the OMD update]
\label{lem:omd-closed-form}
Fix $M_t\in\Delta_H$, $\eta>0$, and $g_t\in\mathbb{R}^H$, and let
$\widetilde{M}_{t+1}$ be defined coordinatewise by \eqref{eq:alg-omd-mult},
i.e.\ $\widetilde{M}_{t+1}^{(i)}=M_t^{(i)}e^{-\eta g_t^{(i)}}$. Then
\[
\begin{aligned}
&\operatorname*{arg\,min}_{M\in\Delta_H}\Bigl\{\eta\langle g_t,M\rangle+D_\Psi(M\|M_t)\Bigr\} \\
&\quad=
\operatorname*{arg\,min}_{M\in\Delta_H} D_\Psi(M\|\widetilde{M}_{t+1})
=
P(\widetilde{M}_{t+1}),
\end{aligned}
\]
where $P(\cdot)$ is the Bregman projection of Lemma~\ref{lem:bregman-projection}.
Consequently, the OMD update \eqref{eq:alg-omd-argmin} coincides exactly with
the closed-form OMD update \eqref{eq:alg-omd-mult}--\eqref{eq:alg-omd-proj},
and in particular $M_{t+1}\in\Delta_H$.
\end{lemma}

\begin{proof}
For any $M\in\mathbb{R}_{\ge0}^H$, expand the objective using the definitions of
$D_\Psi$ from \eqref{eq:def-negentropy} and of $\widetilde{M}_{t+1}$:
\[
\begin{aligned}
&\eta\langle g_t,M\rangle + D_\Psi(M\|M_t) \\
&=\sum_{i=1}^H\Bigl(\eta g_t^{(i)}M^{(i)}
+M^{(i)}\log\tfrac{M^{(i)}}{M_t^{(i)}}-M^{(i)}+M_t^{(i)}\Bigr) \\
&=\sum_{i=1}^H\Bigl(M^{(i)}\log\tfrac{M^{(i)}}{M_t^{(i)}e^{-\eta g_t^{(i)}}}
-M^{(i)}+M_t^{(i)}\Bigr) \\
&=\sum_{i=1}^H\Bigl(M^{(i)}\log\tfrac{M^{(i)}}{\widetilde{M}_{t+1}^{(i)}}
-M^{(i)}+\widetilde{M}_{t+1}^{(i)}\Bigr) \\
&\quad+\sum_{i=1}^H\bigl(M_t^{(i)}-\widetilde{M}_{t+1}^{(i)}\bigr) \\
&=D_\Psi(M\|\widetilde{M}_{t+1})
+\sum_{i=1}^H\bigl(M_t^{(i)}-\widetilde{M}_{t+1}^{(i)}\bigr).
\end{aligned}
\]
The second sum does not depend on $M$, so minimizing the left-hand side over
$M\in\Delta_H$ is equivalent to minimizing $D_\Psi(M\|\widetilde{M}_{t+1})$ over
$M\in\Delta_H$, which by definition is the Bregman projection
$P(\widetilde{M}_{t+1})$. Lemma~\ref{lem:bregman-projection} identifies
$P(\widetilde{M}_{t+1})$ coordinatewise as
$\widetilde{M}_{t+1}^{(i)}/\max(1,S_{t+1})$ with
$S_{t+1}=\sum_j\widetilde{M}_{t+1}^{(j)}$, which is exactly
\eqref{eq:alg-omd-proj}.
\end{proof}

\begin{lemma}[One-step movement of OMD]
\label{lem:omd-move}
For the OMD update with step size \(\eta>0\) and subgradient
\(g_t\in\partial \ell_t(M_t)\),
\[
\|M_{t+1}-M_t\|_1 \le \eta\,\|g_t\|_\infty.
\]
In particular, if \(\|g_t\|_\infty\le L\), then \(\|M_{t+1}-M_t\|_1\le \eta L\).
\end{lemma}

\begin{proof}
Let \(\Psi\) and \(D_\Psi\) be as in \eqref{eq:def-negentropy}. By definition
\eqref{eq:alg-omd-argmin}, \(M_{t+1}\) minimizes the convex function
\(F(M)\triangleq\eta\langle g_t,M\rangle+D_\Psi(M\|M_t)\) over the convex set
\(\Delta_H\), and \(F\) is differentiable on \(\mathbb{R}_{>0}^H\) with
\(\nabla F(M)=\eta g_t+\nabla\Psi(M)-\nabla\Psi(M_t)\). The first-order optimality
condition for constrained convex minimization therefore gives, for all
\(M\in\Delta_H\),
\[ 
\begin{aligned}
\Bigl\langle g_t + \tfrac{1}{\eta}\bigl(\nabla\Psi(M_{t+1})-\nabla\Psi(M_t)\bigr), 
 M-M_{t+1}\Bigr\rangle \ge 0.
\end{aligned}
\]
Taking \(M=M_t\) and rearranging yields
\[
\langle g_t,M_t-M_{t+1}\rangle
\ge
\tfrac{1}{\eta}\Bigl\langle
\nabla\Psi(M_{t+1})-\nabla\Psi(M_t),\,M_{t+1}-M_t\Bigr\rangle.
\]
By the Bregman identity,
\[
\begin{aligned}
&\Bigl\langle \nabla\Psi(M_{t+1})-\nabla\Psi(M_t),\,M_{t+1}-M_t\Bigr\rangle \\
&\qquad=
D_\Psi(M_{t+1}\|M_t)+D_\Psi(M_t\|M_{t+1}).
\end{aligned}
\]
Hence, using Hölder's inequality,
\[
\begin{aligned}
&D_\Psi(M_{t+1}\|M_t)+D_\Psi(M_t\|M_{t+1}) \\
&\qquad\le
\eta\,\|g_t\|_\infty\,\|M_{t+1}-M_t\|_1.
\end{aligned}
\]
Since both iterates lie in \(\Delta_H\), the unnormalized negentropy is
\(1\)-strongly convex with respect to \(\|\cdot\|_1\) on this set: the Hessian
\(\nabla^2\Psi(M)=\mathrm{diag}(1/M^{(i)})\) satisfies
\(v^\top\nabla^2\Psi(M)\,v=\sum_i (v^{(i)})^2/M^{(i)}\ge\|v\|_1^2/(\sum_i
M^{(i)})\ge\|v\|_1^2\) whenever \(\sum_i M^{(i)}\le1\) (Cauchy--Schwarz), so the
usual Pinsker-type bound for the Bregman divergence extends to the sub-probability
simplex \(\Delta_H\),
\begin{align*}
D_\Psi(M_{t+1}\|M_t)&\ge \tfrac12\|M_{t+1}-M_t\|_1^2, \\
D_\Psi(M_t\|M_{t+1})&\ge \tfrac12\|M_{t+1}-M_t\|_1^2,
\end{align*}
and the left-hand side is at least \(\|M_{t+1}-M_t\|_1^2\). Therefore
\[
\begin{aligned}
&\|M_{t+1}-M_t\|_1^2
\le
\eta\,\|g_t\|_\infty\,\|M_{t+1}-M_t\|_1,
\end{aligned}
\]
which implies \(\|M_{t+1}-M_t\|_1\le \eta\|g_t\|_\infty\).
\end{proof}

\begin{lemma}[Cumulative movement]
\label{lem:cum-move}
For every \(t\in\{2,\ldots,T\}\) and every integer \(i\) with \(1\le i\le t-1\),
\[
\|M_{t-i}-M_t\|_1 \le i\,\eta L.
\]
\end{lemma}

\begin{proof}
The range \(1\le i\le t-1\) ensures that every index \(t-j\) for \(j=1,\ldots,i\)
satisfies \(t-j\ge 1\). By the triangle inequality and Lemma~\ref{lem:omd-move},
\[
\|M_{t-i}-M_t\|_1
\le
\sum_{j=1}^i \|M_{t-j}-M_{t-j+1}\|_1
\le i\,\eta L.
\]
\end{proof}

\begin{lemma}[Control-action bounds relative to the benchmark trajectory]
\label{lem:bounding-u-eff}
For all \(t\),
\begin{align*}
u_t^{\mathcal{A}}
&\ge
\langle M_t,\widetilde{w}_t\rangle
- \alpha\,[x_t^{\pi(M_t)}-x_t^{\mathcal{A}}]_+ ,\\[4pt]
u_t^{\mathcal{A}}
&\le
\langle M_t,\widetilde{w}_t\rangle.
\end{align*}
\end{lemma}

\begin{proof}
Recall from \eqref{eq:alg-unconstrained-action}--\eqref{eq:alg-feasible-action} that
\[
\widehat{u}_t = \langle M_t,\widetilde{w}_t\rangle,
\qquad
u_t^{\mathcal{A}}=\min\{\widehat{u}_t,\alpha x_t^{\mathcal{A}}\}.
\]
Therefore,
\[
u_t^{\mathcal{A}}
=
\widehat{u}_t-\bigl[\widehat{u}_t-\alpha x_t^{\mathcal{A}}\bigr]_+
=
\langle M_t,\widetilde{w}_t\rangle
-\bigl[\langle M_t,\widetilde{w}_t\rangle-\alpha x_t^{\mathcal{A}}\bigr]_+.
\]
This identity gives the upper bound immediately, since the subtracted term is
nonnegative.

For the lower bound, Lemma~\ref{lem:appendix-feasibility} applied to the benchmark
policy \(\pi(M_t)\) yields, using \eqref{eq:def-sdac-policy-vector},
\[
u_t^{\pi(M_t)}
=
\langle M_t,\widetilde{w}_t\rangle
\le \alpha x_t^{\pi(M_t)}.
\]
Subtracting \(\alpha x_t^{\mathcal{A}}\) from both sides and applying monotonicity of
\([\cdot]_+\) gives
\[
\begin{aligned}
&\bigl[\langle M_t,\widetilde{w}_t\rangle-\alpha x_t^{\mathcal{A}}\bigr]_+ \\
&\qquad\le
\bigl[\alpha x_t^{\pi(M_t)}-\alpha x_t^{\mathcal{A}}\bigr]_+
=
\alpha\bigl[x_t^{\pi(M_t)}-x_t^{\mathcal{A}}\bigr]_+.
\end{aligned}
\]
Substituting this bound into the expression above for \(u_t^{\mathcal{A}}\) yields
the stated lower bound.
\end{proof}

\begin{lemma}[State-gap recursion]
\label{lem:state-gap-recursion-eff}
Define the state gap
\begin{equation}
\label{eq:def-state-gap}
\Gamma_t^M \triangleq x_t^{\pi(M)} - x_t^{\mathcal{A}}.
\end{equation}
For any fixed $M\in\Delta_H$ and every $t\in\{1,\ldots,T\}$ the following two-sided recursion
holds:
\begin{align}
\Gamma_{t+1}^M
&\ge \alpha\Gamma_t^M + \langle M_t-M,\widetilde{w}_t\rangle - \alpha[\Gamma_t^{M_t}]_+ ,\label{eq:GAMMA-Lower}\\[4pt]
\Gamma_{t+1}^M
&\le \alpha\Gamma_t^M + \langle M_t-M,\widetilde{w}_t\rangle.\label{eq:GAMMA-Upper}
\end{align}
\end{lemma}

\begin{proof}
From the state updates \eqref{eq:sdac-state-update} (for $\pi(M)$) and
\eqref{eq:alg-state-update} (for the algorithm), the arrival $w_t$ cancels and
we obtain the exact identity
\begin{equation}
    \label{eq:state-gap-identity}
    \Gamma_{t+1}^M = x_{t+1}^{\pi(M)}-x_{t+1}^{\mathcal{A}} = \alpha\Gamma_t^M - \bigl(u_t^{\pi(M)}-u_t^{\mathcal{A}}\bigr).
\end{equation}
Consequently, bounding $\Gamma_{t+1}^M$ reduces to bounding the control-action
difference $u_t^{\pi(M)}-u_t^{\mathcal{A}}$.

Using the SDAC representation \eqref{eq:def-sdac-policy-vector} for the benchmark
control action and the bounds of Lemma~\ref{lem:bounding-u-eff} for the algorithm
control action, we have
\[
u_t^{\pi(M)} = \langle M,\widetilde{w}_t\rangle,
\]
and
\[
\langle M_t,\widetilde{w}_t\rangle-\alpha[\Gamma_t^{M_t}]_+ \le u_t^{\mathcal{A}} \le
\langle M_t,\widetilde{w}_t\rangle.
\]
Subtracting the upper bound on $u_t^{\mathcal{A}}$ from $u_t^{\pi(M)}$ yields the
lower bound on $u_t^{\pi(M)}-u_t^{\mathcal{A}}$:
\[
u_t^{\pi(M)}-u_t^{\mathcal{A}}
\ge \langle M,\widetilde{w}_t\rangle - \langle M_t,\widetilde{w}_t\rangle
= - \langle M_t-M,\widetilde{w}_t\rangle.
\]
Inserting this into the identity \eqref{eq:state-gap-identity} gives
\[
\Gamma_{t+1}^M \le \alpha\Gamma_t^M + \langle M_t-M,\widetilde{w}_t\rangle,
\]
which is \eqref{eq:GAMMA-Upper}.

Similarly, using the lower bound on $u_t^{\mathcal{A}}$ produces an upper bound on
$u_t^{\pi(M)}-u_t^{\mathcal{A}}$:
\[
\begin{aligned}
u_t^{\pi(M)}-u_t^{\mathcal{A}}
&\le \langle M,\widetilde{w}_t\rangle - \Bigl(\langle M_t,\widetilde{w}_t\rangle-\alpha[\Gamma_t^{M_t}]_+\Bigr) \\
&= - \langle M_t-M,\widetilde{w}_t\rangle + \alpha[\Gamma_t^{M_t}]_+.
\end{aligned}
\]
Substituting this bound into the identity \eqref{eq:state-gap-identity} yields
\[
\Gamma_{t+1}^M \ge \alpha\Gamma_t^M + \langle M_t-M,\widetilde{w}_t\rangle - \alpha[\Gamma_t^{M_t}]_+,
\]
which is \eqref{eq:GAMMA-Lower} and completes the proof.
\end{proof}

\begin{lemma}[State deviation bound]
\label{lem:state-gap-bound}
For every $t\in\{1,\ldots,T\}$, let the state gap be defined by \eqref{eq:def-state-gap}. Then
\[
-\frac{\alpha\,\eta L}{(1-\alpha)^3}
\le
\Gamma_t^{M_t}
=
x_t^{\pi(M_t)}-x_t^{\mathcal{A}}
\le
\frac{\alpha\,\eta L}{(1-\alpha)^2}.
\]
\end{lemma}

\begin{proof}
From \eqref{eq:GAMMA-Upper}, for any fixed benchmark $M\in\Delta_H$ and every
$s\in\{1,\ldots,T\}$,
\[
\Gamma_{s+1}^{M}
\le
\alpha\Gamma_s^{M} + \langle M_s-M,\widetilde{w}_s\rangle.
\]
Applying the recursion repeatedly up to time $t$ yields
\[
\begin{aligned}
\Gamma_t^{M}
&\le \alpha\Gamma_{t-1}^{M} + \langle M_{t-1}-M,\widetilde{w}_{t-1}\rangle \\
&\le \alpha^2\Gamma_{t-2}^{M} + \alpha\langle M_{t-2}-M,\widetilde{w}_{t-2}\rangle \\
&\qquad\qquad + \langle M_{t-1}-M,\widetilde{w}_{t-1}\rangle \\
&\ \vdots \\
&\le \alpha^{t-1}\Gamma_1^{M}
 + \sum_{s=1}^{t-1}\alpha^{t-1-s}\langle M_s-M,\widetilde{w}_s\rangle.
\end{aligned}
\]
Now set the benchmark equal to the online parameter $M=M_t$. 
Since $x_1^{\pi(M_t)}=x_1^{\mathcal{A}}=0$, we have $\Gamma_1^{M_t}=0$, and therefore
\[
\Gamma_t^{M_t}
\le
\sum_{s=1}^{t-1}\alpha^{t-1-s}\langle M_s-M_t,\widetilde{w}_s\rangle.
\]
Note that $\|\widetilde{w}_s\|_\infty = \max_{i\le H}\alpha^i w_{s-i} \le \alpha$
since $w_\tau\in[0,1]$ and $i\ge1$. Hence,
\[
\begin{aligned}
\Gamma_t^{M_t}
&\le \sum_{s=1}^{t-1}\alpha^{t-1-s}\bigl|\langle M_s-M_t,\widetilde{w}_s\rangle\bigr| \\
&\le \sum_{s=1}^{t-1}\alpha^{t-1-s}\|M_s-M_t\|_1\|\widetilde{w}_s\|_\infty \\
&\le \alpha\,\eta L\sum_{s=1}^{t-1}\alpha^{t-1-s}(t-s) \\
&= \alpha\,\eta L\sum_{j=1}^{t-1} j\,\alpha^{j-1}
\le \alpha\,\eta L\sum_{j=1}^{\infty} j\,\alpha^{j-1}
= \frac{\alpha\,\eta L}{(1-\alpha)^2}.
\end{aligned}
\]
Where we used the geometric-series identity
\[
\sum_{j=1}^{\infty} j\,\alpha^{j-1} = \frac{1}{(1-\alpha)^2}.
\]

Similarly, for the lower bound, unrolling \eqref{eq:GAMMA-Lower} with the same fixed
benchmark $M=M_t$ yields
\[
\Gamma_t^{M_t}
\ge
\sum_{s=1}^{t-1}\alpha^{t-1-s}\langle M_s-M_t,\widetilde{w}_s\rangle
-
\sum_{s=1}^{t-1}\alpha^{t-s}[\Gamma_s^{M_s}]_+.
\]
The upper bound just proved applies to every earlier time $s$, so
$[\Gamma_s^{M_s}]_+\le \alpha\eta L/(1-\alpha)^2$. Also, by Hölder's inequality and
Lemma~\ref{lem:cum-move},
\begin{align*}
\Gamma_t^{M_t}
&\ge
-\alpha\,\eta L\sum_{s=1}^{t-1}\alpha^{t-1-s}(t-s) \\
&\quad -\frac{\alpha\,\eta L}{(1-\alpha)^2}\sum_{s=1}^{t-1}\alpha^{t-s} \\
&= -\alpha\,\eta L\sum_{j=1}^{t-1} j\,\alpha^{j-1}
 -\frac{\alpha\,\eta L}{(1-\alpha)^2}\sum_{j=1}^{t-1}\alpha^j.
\end{align*}
Extending both sums to infinity gives
\[
\Gamma_t^{M_t}
\ge
-\alpha\,\eta L\sum_{j=1}^{\infty} j\,\alpha^{j-1}
-\frac{\alpha\,\eta L}{(1-\alpha)^2}\sum_{j=1}^{\infty}\alpha^j.
\]
Using the geometric-series identities
\[
\sum_{j=1}^{\infty} j\,\alpha^{j-1} = \frac{1}{(1-\alpha)^2},
\qquad
\sum_{j=1}^{\infty}\alpha^j = \frac{\alpha}{1-\alpha},
\]
we obtain
\[
\Gamma_t^{M_t}
\ge
-\frac{\alpha\,\eta L}{(1-\alpha)^2}
-\frac{\alpha^2\,\eta L}{(1-\alpha)^3}
=
-\frac{\alpha\,\eta L}{(1-\alpha)^3}.
\]
Hence, the claim.
\end{proof}

\begin{lemma}[Action-state gap]
\label{lem:u-gap-by-b-gap-eff}
Let the state gap be defined by \eqref{eq:def-state-gap}. Then for all \(t\in\{1,\ldots,T\}\),
\[
0
\le
u_t^{\pi(M_t)}-u_t^{\mathcal{A}}
\le
\frac{\alpha^2\,\eta L}{(1-\alpha)^2}.
\]
\end{lemma}

\begin{proof}
By \eqref{eq:def-sdac-policy-vector} and \eqref{eq:alg-unconstrained-action}, the
benchmark control action equals the unconstrained control action, \(u_t^{\pi(M_t)} =
\langle M_t,\widetilde{w}_t\rangle = \widehat{u}_t\), while
\(u_t^{\mathcal{A}}=\min\{\widehat{u}_t,\alpha x_t^{\mathcal{A}}\}\le \widehat{u}_t\)
by \eqref{eq:alg-feasible-action}. Hence the gap is nonnegative and
\[
u_t^{\pi(M_t)} - u_t^{\mathcal{A}}
=
\widehat{u}_t - \min\{\widehat{u}_t,\alpha x_t^{\mathcal{A}}\}
=
\bigl[\langle M_t,\widetilde{w}_t\rangle - \alpha x_t^{\mathcal{A}}\bigr]_+ .
\]
As established in the proof of Lemma~\ref{lem:bounding-u-eff} (using feasibility of
\(\pi(M_t)\)),
\[
\bigl[\langle M_t,\widetilde{w}_t\rangle - \alpha x_t^{\mathcal{A}}\bigr]_+
\le
\alpha\,[x_t^{\pi(M_t)}-x_t^{\mathcal{A}}]_+
=
\alpha\,[\Gamma_t^{M_t}]_+ .
\]
By Lemma~\ref{lem:state-gap-bound},
\([\Gamma_t^{M_t}]_+ \le \dfrac{\alpha\,\eta L}{(1-\alpha)^2}\), so
\[
u_t^{\pi(M_t)}-u_t^{\mathcal{A}}
\le
\alpha\,[\Gamma_t^{M_t}]_+
\le
\frac{\alpha^2\,\eta L}{(1-\alpha)^2},
\]
which completes the proof.
\end{proof}

\begin{lemma}[Stability term]
\label{lem:stab-sum}
For all horizons \(T\),
\[
\begin{aligned}
\sum_{t=1}^T \Bigl(c_t(x_t^{\mathcal{A}},u_t^{\mathcal{A}})
-c_t\bigl(x_t^{\pi(M_t)},u_t^{\pi(M_t)}\bigr)\Bigr) 
\le
\frac{\alpha\,T\eta L^2}{(1-\alpha)^2},
\end{aligned}
\]
where \(L=\alpha L_u+\tfrac{L_x}{1-\alpha}\).
\end{lemma}

\begin{proof}
By Lipschitz continuity of \(c_t\) (see equation \eqref{eq:model-lipschitz}),
\[
\begin{aligned}
&c_t\bigl(x_t^{\mathcal{A}},u_t^{\mathcal{A}}\bigr)
-
c_t\bigl(x_t^{\pi(M_t)},u_t^{\pi(M_t)}\bigr) \\
&\qquad\le
L_u \big|u_t^{\mathcal{A}} - u_t^{\pi(M_t)}\big|
+ L_x \big|x_t^{\mathcal{A}} - x_t^{\pi(M_t)}\big|.
\end{aligned}
\]
From Lemma~\ref{lem:u-gap-by-b-gap-eff},
\[
|u_t^{\mathcal{A}} - u_t^{\pi(M_t)}|
\le \frac{\alpha^2\,\eta L}{(1-\alpha)^2},
\]
Also, by Lemma~\ref{lem:state-gap-bound} and \eqref{eq:def-state-gap},
\[
|x_t^{\mathcal{A}} - x_t^{\pi(M_t)}|
= |\Gamma_t^{M_t}|
\le \frac{\alpha\,\eta L}{(1-\alpha)^3}.
\]
Therefore, for each \(t\in\{1,\ldots,T\}\),
\[
\begin{aligned}
&c_t\bigl(x_t^{\mathcal{A}},u_t^{\mathcal{A}}\bigr)
-
c_t\bigl(x_t^{\pi(M_t)},u_t^{\pi(M_t)}\bigr) \\
&\qquad\le
L_u \frac{\alpha^2\,\eta L}{(1-\alpha)^2}
+ L_x \frac{\alpha\,\eta L}{(1-\alpha)^3}.
\end{aligned}
\]
Summing over \(t=1,2,\dots,T\) gives
\[
\begin{aligned}
&\sum_{t=1}^T
\Bigl(
c_t(x_t^{\mathcal{A}},u_t^{\mathcal{A}})
-
c_t(x_t^{\pi(M_t)},u_t^{\pi(M_t)})
\Bigr) \\
&\qquad\le
\frac{\alpha\,T\eta L}{(1-\alpha)^2}\left(
\alpha L_u 
+  \frac{L_x}{1-\alpha}
\right)
=
\frac{\alpha\,T\eta L^2}{(1-\alpha)^2},
\end{aligned}
\]
where we used \(L = \alpha L_u + \tfrac{L_x}{1-\alpha}\) in the final step.
\end{proof}



\begin{proof}[Proof of Theorem~\ref{thm:regret-bound}]

Starting from \eqref{eq:def-regret} we have the following decomposition of the
regret:
\begin{align*}
\RegSDAC{T}
&= \sum_{t=1}^T c_t(x_t^{\mathcal{A}},u_t^{\mathcal{A}}) \\
&\quad - \min_{M\in\Delta_H}\sum_{t=1}^T c_t\bigl(x_t^{\pi(M)},u_t^{\pi(M)}\bigr) \\
&= \sum_{t=1}^T c_t(x_t^{\mathcal{A}},u_t^{\mathcal{A}})
 - \sum_{t=1}^T c_t\bigl(x_t^{\pi(M_t)},u_t^{\pi(M_t)}\bigr) \\
&\quad + \sum_{t=1}^T c_t\bigl(x_t^{\pi(M_t)},u_t^{\pi(M_t)}\bigr) \\
&\quad - \min_{M\in\Delta_H}\sum_{t=1}^T c_t\bigl(x_t^{\pi(M)},u_t^{\pi(M)}\bigr).
\end{align*} 
We call the first sum the stability term, and the second sum the OMD term. We will
bound each term separately.

We bound the OMD term using the OMD update \eqref{eq:alg-omd-argmin} with
the unnormalized negentropy $\Psi$ and Bregman divergence $D_\Psi$ from
\eqref{eq:def-negentropy}. Under the initialization \eqref{eq:alg-init}, for any
benchmark $M^\star\in\Delta_H$ with $W^\star\triangleq\sum_{i=1}^H M^{\star(i)}\le
1$,
\[
\begin{aligned}
&D_\Psi(M^\star\|M_1) \\
&\quad=
\sum_{i=1}^H\Bigl(M^{\star(i)}\log\bigl((H+1)M^{\star(i)}\bigr)
-M^{\star(i)}+\tfrac{1}{H+1}\Bigr) \\
&\quad=
\sum_{i=1}^H M^{\star(i)}\log M^{\star(i)}
+
W^\star\log(H+1)
-
W^\star
+
\tfrac{H}{H+1}.
\end{aligned}
\]
Since each $M^{\star(i)}\in[0,1]$, one has $M^{\star(i)}\log M^{\star(i)}\le 0$, and
therefore
\[
D_\Psi(M^\star\|M_1)
\le
W^\star\log(H+1)-W^\star+\tfrac{H}{H+1}.
\]
The right-hand side is at most $\log(H+1)$ for every $W^\star\in[0,1]$: equivalently,
\[
(1-W^\star)\log(H+1)+W^\star-\tfrac{H}{H+1}\ge 0,
\]
which holds because the left-hand side is affine in $W^\star$ and nonnegative at the
endpoints $W^\star=0$ and $W^\star=1$. Hence $D_\Psi(M^\star\|M_1)\le\log(H+1)$. The
surrogate costs $\ell_t$ from \eqref{eq:alg-surrogate-cost} are convex and
$L$-Lipschitz by Lemma~\ref{lem:appendix-surrogate}, so the standard OMD regret bound
with Bregman divergence $D_\Psi$ on $\Delta_H$ \cite{hazan2016introduction} yields
\begin{align*}
\sum_{t=1}^T \ell_t(M_t) - \min_{M\in\Delta_H}\sum_{t=1}^T \ell_t(M)
&\le
\frac{\log(H+1)}{\eta} + \frac{\eta L^2 T}{2}.
\end{align*}
Using \eqref{eq:alg-surrogate-cost}, we substitute
\(\ell_t(M)=c_t\bigl(x_t^{\pi(M)},u_t^{\pi(M)}\bigr)\), so that
\begin{align*}
\sum_{t=1}^T \ell_t(M_t)
&=
\sum_{t=1}^T c_t\bigl(x_t^{\pi(M_t)},u_t^{\pi(M_t)}\bigr), \\
\min_{M\in\Delta_H}\sum_{t=1}^T \ell_t(M)
&=
\min_{M\in\Delta_H}\sum_{t=1}^T c_t\bigl(x_t^{\pi(M)},u_t^{\pi(M)}\bigr).
\end{align*}
Thus this term becomes
\begin{align*}
&\sum_{t=1}^T c_t\bigl(x_t^{\pi(M_t)},u_t^{\pi(M_t)}\bigr) \\
&\quad-
\min_{M\in\Delta_H}\sum_{t=1}^T c_t\bigl(x_t^{\pi(M)},u_t^{\pi(M)}\bigr) \\
&\qquad\le
\frac{\log(H+1)}{\eta} + \frac{\eta L^2 T}{2}.
\end{align*}


From Lemma~\ref{lem:stab-sum}, the stability term satisfies
\[
\begin{aligned}
&\sum_{t=1}^T c_t(x_t^{\mathcal{A}},u_t^{\mathcal{A}})
- \sum_{t=1}^T c_t(x_t^{\pi(M_t)},u_t^{\pi(M_t)}) \\
&\qquad\le 
\frac{\alpha\,\eta L^2 T}{(1-\alpha)^2}.
\end{aligned}
\]
Adding the OMD term and the stability term yields the combined bound
\[
\begin{aligned}
\RegSDAC{T}
\;\le\;
\frac{\log(H+1)}{\eta}
\;+\;
\frac{\eta L^2 T}{2}
\;+\;
\frac{\alpha\,\eta L^2 T}{(1-\alpha)^2}.
\end{aligned}
\]
Rearranging by factoring out the \(\eta\) term gives
\[
\RegSDAC{T}
\;\le\;
\frac{\log(H+1)}{\eta}
\;+\;
\eta L^2 T\Bigl(\tfrac{1}{2}+\tfrac{\alpha}{(1-\alpha)^2}\Bigr).
\]
To optimize over \(\eta > 0\), we apply the inequality \(A/\eta + B\eta \le
2\sqrt{AB}\) with \(A=\log(H+1)\) and \(B= L^2
T(\tfrac{1}{2}+\tfrac{\alpha}{(1-\alpha)^2})\). This yields the optimal step size
\[
\eta^* = \sqrt{\frac{\log(H+1)}{L^2
T\bigl(\tfrac{1}{2}+\tfrac{\alpha}{(1-\alpha)^2}\bigr)}}
\] 
and the resulting regret bound
\[
\RegSDAC{T} \;\le\;
2L\sqrt{\Bigl(\tfrac{1}{2}+\tfrac{\alpha}{(1-\alpha)^2}\Bigr)T\log(H+1)}.
\]
\end{proof}



\bibliographystyle{IEEEtran}
\bibliography{ref}

\end{document}


